\section{The \OurLang\ Calculus: Syntax, Semantics and Types}
\label{sec:calculus}

\subsection{Syntax and Semantics}

\Cref{fig:syntax} gives the syntax of \OurLang, a concurrent \lcalc\ with
synchronous message passing, loosely based on the formulation of GV by
\citet{conf/esop/LindleyM15}.
Variables are ranged over by $\vr, \vrB$ and channels by $\ch, \chB$.
$\tApp{\tNew}{\tUnit}$ creates a channel and returns both of its ends,
$\tApp{\tSpawn}{\val}$ runs $\tApp{\val}{\tUnit}$ in a new thread,
$\tApp{\tRecv}{\ch}$ receives a message on $\ch$ and returns it with the
continuation channel,
$\tApp{\tSend}{\tPair{\val}{\ch}}$ sends $\val$ on $\ch$,
$\tApp{\tClos}{\ch}$ closes $\ch$,
$\tApp{\tLink}{\tPair{\ch}{\chB}}$ forwards between two ends, and
$\tApp{\tForce}{\tBoxP{\vr}{\tm}}$ unrolls a box.
A box $\tBoxP{\vr}{\tm}$ binds $\vr$ to itself in $\tm$, which is the
only form of recursion in the calculus.

\begin{figure}[thp]
  \adjustfigure[]
  \rulesection*{Constants $\cst$, values $\val$, terms $\tm$ (\rocq{expr})}
  \begin{align*}
    \cst &::=
      \tNew \sepr \tSpawn \sepr \tRecv \sepr \tSend \sepr
      \tClos \sepr \tLink \sepr \tForce
    & \text{(constants)}
    \\
    \val &::=
      \tUnit \sepr \cst \sepr \ch \sepr \tAbs{\vr}{\tm} \sepr
      \tBoxP{\vr}{\tm} \sepr \tPair{\val}{\val} \sepr \tInj{i}{\val}
    & \text{(values)}
    \\
    \tm &::=
      \vr \sepr \val \sepr \tApp{\tm}{\tm} \sepr \tPair{\tm}{\tm} \sepr
      \tLet*{\tm}{\tm} \sepr \tLet[\vr][\vrB]{\tm}{\tm}
    \\ & \qquad {}\sepr
      \tInj{i}{\tm} \sepr \tCase{\tm}{\ivx}{\tm}{\ivy}{\tm}
    & \text{(terms)}
  \end{align*}
  \rulesection{Derived forms}
  \begin{align*}
    \tSeq{\tm_1}{\tm_2} &:= \tLet*{\tm_1}{\tm_2}
    &
    \tAbs{\tUnit}{\tm} &:= \tAbs{\ivx}{\tLet*{\ivx}{\tm}}
    &
    \tLetS{\vr}{\tm_1}{\tm_2} &:= \tApp{(\tAbs{\vr}{\tm_2})}{\tm_1}
  \end{align*}
  \rulesectionend
  \caption{Syntax of \OurLang (\rocq{expr})}
  \Description{Grammars for the constants, values and terms of GV-sharp,
    and three derived forms.}
  \label{fig:syntax}
\end{figure}

\begin{figure}[tp]
  \adjustfigure[]
  \rulesection*{Session labels $\sigma$, term labels $\lambda$, evaluation contexts $\ectx$ (\rocq{slabel}, \rocq{tlabel})}
  \begin{align*}
    \sigma &::=
      \lRecv{\ch}{\val} \sepr
      \lSend{\ch}{\val} \sepr
      \lTerm{\ch} \sepr
      \lLink{\ch}{\chB}
    & \text{(session labels)}
    \\
    \lambda &::=
      \beta \sepr
      \sigma \sepr
      \lNew{\ch} \sepr
      \lSpawnL{\val}
    & \text{(term labels)}
    \\
    \ectx &::=
      \hole \sepr
      \tApp{\ectx}{\tm} \sepr
      \tApp{\val}{\ectx} \sepr
      \tPair{\ectx}{\tm} \sepr
      \tPair{\val}{\ectx} \sepr
      \tLet*{\ectx}{\tm} \sepr
      \tLet[\vr][\vrB]{\ectx}{\tm}
    \\ & \qquad {}\sepr
      \tInj{1}{\ectx} \sepr \tInj{2}{\ectx} \sepr
      \tCase{\ectx}{\ivx_1}{\tm_1}{\ivx_2}{\tm_2}
    & \text{(eval.\ contexts)}
  \end{align*}
  \rulesection{Labelled reductions for terms (\rocq{tstep})}
  \begin{proofrules}
    \infer*[lab=T-Act-App]{}{
      \tApp{(\tAbs{\vr}{\tm})}{\val} \trans{\beta} \tSub{\tm}{\val}{\vr}
    } \label{rule:T-Act-App}
    \and
    \infer*[lab=T-Act-Let*]{}{
      \tLet*{\tUnit}{\tm} \trans{\beta} \tm
    } \label{rule:T-Act-Let*}
    \and
    \infer*[lab=T-Act-LetPair]{}{
      \tLet[\vr][\vrB]{\tPair{\valB}{\val}}{\tm} \trans{\beta} \tSub{\tSub{\tm}{\valB}{\vr}}{\val}{\vrB}
    } \label{rule:T-Act-LetPair}
    \and
    \infer*[lab=T-Act-Case]{}{
      \tCase{(\tInj{\asTm{i}}{\val})}{\ivx_1}{\tm_1}{\ivx_2}{\tm_2} \trans{\beta} \tSub{\tm_{\asTm{i}}}{\val}{\ivx_{\asTm{i}}}
    } \label{rule:T-Act-Case}
    \and
    \infer*[lab=T-Act-Force]{}{
      \tApp{\tForce}{\tBoxP{\vr}{\tm}} \trans{\beta} \tSub{\tm}{\tBoxP{\vr}{\tm}}{\vr}
    } \label{rule:T-Act-Force}
    \and
    \infer*[lab=T-Act-New]{}{
      \tApp{\tNew}{\tUnit} \trans{\lNew{\ch}} \tPair{\ch}{\ch}
    } \label{rule:T-Act-New}
    \and
    \infer*[lab=T-Act-Spawn]{}{
      \tApp{\tSpawn}{\val} \trans{\lSpawnL{\val}} \tUnit
    } \label{rule:T-Act-Spawn}
    \and
    \infer*[lab=T-Act-Recv]{}{
      \tApp{\tRecv}{\ch} \trans{\lRecv{\ch}{\val}} \tPair{\val}{\ch}
    } \label{rule:T-Act-Recv}
    \and
    \infer*[lab=T-Act-Send]{}{
      \tApp{\tSend}{\tPair{\val}{\ch}} \trans{\lSend{\ch}{\val}} \ch
    } \label{rule:T-Act-Send}
    \and
    \infer*[lab=T-Act-Close]{}{
      \tApp{\tClos}{\ch} \trans{\lTerm{\ch}} \tUnit
    } \label{rule:T-Act-Close}
    \and
    \infer*[lab=T-Act-Link]{}{
      \tApp{\tLink}{\tPair{\ch}{\chB}} \trans{\lLink{\ch}{\chB}} \tUnit
    } \label{rule:T-Act-Link}
    \and
    \infer*[lab=T-Act-Ctx]{
      \tm \trans{\lambda} \tm'
    }{
      \ectx[\tm] \trans{\lambda} \ectx[\tm']
    } \label{rule:T-Act-Ctx}
  \end{proofrules}
  \rulesectionend
  \caption{Semantics of \OurLang\ terms}
  \Description{Grammars for evaluation contexts, session labels and term
    labels, and the eleven axioms plus one congruence rule of the
    labelled term reduction relation.}
  \label{fig:term-semantics}
\end{figure}

Both ends of a channel carry one name: \ref{rule:T-Act-New} returns the
pair $\tPair{\ch}{\ch}$. The subjects of a label are
$\fc(\beta) = \fc(\lNew{\ch}) = \fc(\lSpawnL{\val}) = \emptyset$,
$\fc(\lRecv{\ch}{\val}) = \fc(\lSend{\ch}{\val}) = \fc(\lTerm{\ch}) = \set{\ch}$
and $\fc(\lLink{\ch}{\chB}) = \set{\ch,\chB}$, and the channels a label
sends away are $\sent(\lSend{\ch}{\val}) = \fchans(\val)$ and
$\emptyset$ otherwise.
A forwarding label carries an unordered pair, so
$\lLink{\ch}{\chB} = \lLink{\chB}{\ch}$, and
$\tApp{\tLink}{\tPair{\ch}{\ch}}$ emits the singleton
$\lLinkSingle{\ch}$.
We write $\tm \redu \tm'$ for $\tm \trans{\beta} \tm'$ and $\redu[*]$ for
its reflexive transitive closure.


A \emph{process} $\proc$ is a pair $(\chset, \pool)$ of a finite set of
restricted channels $\chset$ and a finite multiset of closed terms
$\pool$ running in parallel.
We write $S_1 \disjoint S_2$ and $S_1 \cupdot S_2$ for disjointness and
disjoint union of sets, and $\pool_1 \uplus \pool_2$ for multiset union.
\Cref{fig:proc-semantics} defines the labelled transition relation
$(\chset,\pool) \trans{\pi} (\chset,\pool)$.
A process may represent a whole network or a sub-network, in which case
the terms of $\pool$ use channels outside $\chset$ to talk to the
environment.
\Cref{rule:P-Act-Session} lifts every offer on a channel shared with the
environment, forwarding included, and removes the sent channels from
$\chset$ in the spirit of the \textsc{Open} rule of the standard LTS for
the \picalc~\cite{lts-pi}.

\begin{remark}
\label{rm:sentc}
Computing the sent channels \emph{syntactically} is not a neutral
choice. In the pure \picalc\ every message is a channel name, while in
\OurLang\ a message is an arbitrary value. Sending
$\val = (\tAbs{\vr}{\asTm{\Omega}\ \ch})$ with $\asTm{\Omega}$ diverging
does not share $\ch$, and still $\ch \in \sent(\val)$, so the LTS
removes it from the local scope. This sound approximation buys a simpler
LTS at the price of two syntactic side conditions in the logical
relation (\cref{rm:vabs}).
\end{remark}

\begin{theorem}[Harmony \rocq{harmony}]
  \label{thm:harmony}
  $\proc \trans{\tau} \proc'$
  if and only if
  $\proc \sred \proc'$
  (\cref{sec:app-harmony}, \cref{fig:red-semantics}).
\end{theorem}

\begin{proof}
  \Cref{sec:app-harmony}, after the reduction semantics it compares the
  LTS with.
\end{proof}

Transitions $\proc \trans{\tau} \proc'$ are written $\proc \redu \proc'$
and called process reductions, with $\trans{}[*]$ for the reflexive
transitive closure.

\begin{figure}[tp]
  \adjustfigure[]
  \rulesection{Labelled reductions for processes (\rocq{plabel})}
  \begin{align*}
    \pi &::= \tau \sepr \sigma
    \qquad \text{(process labels)}
  \end{align*}
  \begin{proofrules}
    \infer*[lab=P-Act-Session]{
      \tm \trans{\sigma} \tm' \\
      \fc(\sigma) \disjoint \chset
    }{
      (\chset, \pool \uplus \mset{\tm}) \trans{\sigma} (\chset \setminus \sent(\sigma), \pool \uplus \mset{\tm'})
    } \label{rule:P-Act-Session}
    \and
    \infer*[lab=P-Act-Beta]{
      \tm \trans{\beta} \tm'
    }{
      (\chset, \pool \uplus \mset{\tm}) \trans{\tau} (\chset, \pool \uplus \mset{\tm'})
    } \label{rule:P-Act-Beta}
    \and
    \infer*[lab=P-Act-New]{
      \ch~\text{fresh}
      \\
      \smash{\tm \trans{\lNew{\ch}} \tm'}
    }{
      (\chset, \pool \uplus \mset{\tm}) \trans{\tau} (\chset \cupdot \set{\ch}, \pool \uplus \mset{\tm'})
    } \label{rule:P-Act-New}
    \and
    \infer*[lab=P-Act-Spawn]{
      \smash{\tm \trans{\lSpawnL{\val}} \tm'}
    }{
      (\chset, \pool \uplus \mset{\tm}) \trans{\tau} (\chset, \pool \uplus \mset{\tm', \tApp{\val}{\tUnit}})
    } \label{rule:P-Act-Spawn}
    \and
    \infer*[lab=P-Act-Msg]{
      \tm \trans{\lSend{\ch}{\val}} \tm' \\
      \tmB \trans{\lRecv{\ch}{\val}} \tmB'
    }{
      (\chset, \pool \uplus \mset{\tm, \tmB}) \trans{\tau} (\chset, \pool \uplus \mset{\tm', \tmB'})
    } \label{rule:P-Act-Msg}
    \and
    \infer*[lab=P-Act-Close]{
      \tm \trans{\lTerm{\ch}} \tm' \\
      \tmB \trans{\lTerm{\ch}} \tmB' \\
      \ch \in \chset
    }{
      (\chset, \pool \uplus \mset{\tm, \tmB}) \trans{\tau} (\chset \setminus \set{\ch}, \pool \uplus \mset{\tm', \tmB'})
    } \label{rule:P-Act-Close}
    \and
    \infer*[lab=P-Act-Link]{
      \tm \trans{\lLink{\ch}{\chB}} \tm' \\
      \ch \in \chset
    }{
      (\chset, \pool \uplus \mset{\tm}) \trans{\tau} (\chset \setminus \set{\ch}, \tSub{(\pool \uplus \mset{\tm'})}{\chB}{\ch})
    } \label{rule:P-Act-Link}
  \end{proofrules}
  \rulesectionend
  \caption{LTS semantics of \OurLang\ processes}
  \Description{The grammar of process labels and the seven rules of the
    process labelled transition relation: one lifting an external offer,
    and one each for beta steps, channel creation, spawning, matched
    messages, matched closes and forwarding.}
  \label{fig:proc-semantics}
\end{figure}

\paragraph{Properties of processes}
Forked threads are never joined, so they are expected to reduce to unit,
while the initial thread may evaluate to the value computed by the whole
program.\footnote{This design decision is from
  \citet{conf/esop/LindleyM15}. The rest of the development does
  \emph{not} depend on it.}

\begin{definition}[Irreducible terms and processes \rocq{irred}]
We write $\irred \tm$ if there are no $\lambda, \tm'$ with
$\tm \trans{\lambda} \tm'$, and $\irred \proc$ if there are no
$\pi, \proc'$ with $\proc \trans{\pi} \proc'$.
\end{definition}

\begin{definition}[Terminated process \rocq{terminated}]
\label{def:terminated}
  A process $\proc$ is \emph{terminated} if
  $\proc = (\asPrOp{\emptyset}, \mset{\val,\tUnit,\dots,\tUnit})$
  for some~$\val$.
\end{definition}

\begin{definition}[Safe process \rocq{safe}]
\label{def:safe}
  A process $\proc$ is \emph{safe} if it is terminated, or if
  $\proc \redu \proc'$ for some~$\proc'$.
\end{definition}

\begin{definition}[Deadlock freedom \rocq{deadlock_free}]
  \label{def:df}
  A process $\proc$ is \emph{deadlock free} if
  ${\forall \proc'.\,\proc \redu[*] \proc'}$ implies $\proc'$ is safe.
\end{definition}

\begin{example}[An exchange between two threads \rocq{sync_main}]
\label{ex:sync}
  \begin{align*}
    \asTmId{sync} := {}
      & \tLetS{\tPair{\ch}{\chB}}{\tApp{\tNew}{\tUnit}}{}
        \tSeq{\tApp{\tSpawn}{(\tAbs{\tUnit}{\tApp{\tClos}{(\tApp{\tSend}{\tPair{\tUnit}{\ch}})}})}}{}
      \\
      & \tLet[\ivx][\chB]{\tApp{\tRecv}{\chB}}{\tSeq{\ivx}{\tApp{\tClos}{\chB}}}
  \end{align*}
  The spawned thread sends $\tUnit$ on $\ch$ and closes it, the main
  thread receives on $\chB$ and closes it. The process
  $(\emptyset, \mset{\asTmId{sync}})$ has exactly one reduction sequence,
  and it ends terminated, so it is safe and deadlock free
  (\rocq{sync_deadlock_free}).
\end{example}

\begin{example}[A cyclic topology \rocq{cycle_main}]
\label{ex:cyclic}
  \begin{align*}
    \asTmId{cycle} := {}
      & \tLetS{\tPair{\ch}{\ch'}}{\tApp{\tNew}{\tUnit}}{}
        \tLetS{\tPair{\chB}{\chB'}}{\tApp{\tNew}{\tUnit}}{}
      \\
      & \tSeq{\tApp{\tSpawn}{(\tAbs{\tUnit}{\tSeq{\tApp{\tClos}{\ch'}}{\tApp{\tClos}{\chB'}}})}}{}
        \tSeq{\tApp{\tClos}{\ch}}{\tApp{\tClos}{\chB}}
  \end{align*}
  Two threads share the two channels $\ch$ and $\chB$, so the topology
  is cyclic. Both threads close $\ch$ before $\chB$, the process has
  exactly one reduction sequence, and it is deadlock free
  (\rocq{cycle_deadlock_free}).
\end{example}

\begin{example}[A deadlocked process \rocq{blocked_main}]
\label{ex:deadlock}
  Reversing the two closes of one thread of \cref{ex:cyclic} gives
  \begin{align*}
    \asTmId{blocked} := {}
      & \tLetS{\tPair{\ch}{\ch'}}{\tApp{\tNew}{\tUnit}}{}
        \tLetS{\tPair{\chB}{\chB'}}{\tApp{\tNew}{\tUnit}}{}
      \\
      & \tSeq{\tApp{\tSpawn}{(\tAbs{\tUnit}{\tSeq{\tApp{\tClos}{\chB'}}{\tApp{\tClos}{\ch'}}})}}{}
        \tSeq{\tApp{\tClos}{\ch}}{\tApp{\tClos}{\chB}}
  \end{align*}
  After the two creations and the spawn the process is irreducible and
  not terminated: one thread waits to close $\ch$, the other waits to
  close $\chB$. It is not safe, so no process that reaches it is
  deadlock free.
\end{example}

\subsection{The Static System}
\label{sec:syntacticTyping}
\label{sec:syntacticTyping:typeSys}

Types carry \emph{points}. A point is a name with a natural offset, a
channel type is $\sAt{\sty}{\pntA}$, and the $k$-th action of a channel
at $\pntA$ happens at $\poff{\pntA}{k}$. \Cref{fig:points} gives points,
facts and constraints, \cref{fig:types} the types, and
\cref{fig:facts} the fact theory the rules read every side condition
from.

\begin{figure}[tp]
  \adjustfigure[]
  \rulesection*{Names $\pnm$, points $\pntA$, facts $\fact$, constraints $\cstr$ (\rocq{point}, \rocq{fact}, \rocq{bcs})}
  \begin{align*}
    \pnm &::= \pcst \sepr \pviA \sepr \pslot{\pntA}{j}
    & \text{(a constant, a point variable, the $j$-th channel carried at $\pntA$)}
    \\
    \pntA &::= \poff{\pnm}{k}
    & \text{(the $k$-th action of $\pnm$, $k \in \Nat$; $\poff{\pnm}{0}$ is written $\pnm$)}
    \\
    \fact \in \tyFact &::= \fLt{\pntA}{\pntB} \sepr \fEq{\pntA}{\pntB}
    & \text{($\pntA$ acts before $\pntB$; $\pntA$ and $\pntB$ are one point)}
    \\
    \cstr &\subseteq \tyFact
    & \text{(a set of facts, finite or infinite, under a binder $\pvivec$)}
  \end{align*}
  \rulesection{Roots}
  \begin{align*}
    \rootof(\poff{\pviA}{k}) &= \pviA
    &
    \rootof(\poff{\pslot{\pntA}{j}}{k}) &= \rootof(\pntA)
    &
    \rootof(\poff{\pcst}{k}) &= \pcst
  \end{align*}
  \rulesectionend
  \caption{Points, facts and constraints}
  \Description{The grammars of names, points, facts and constraints, and
    the three equations defining the root of a point.}
  \label{fig:points}
\end{figure}

A channel at $\pntA$ acts there and moves to $\poff{\pntA}{1}$. A
message sent at $\pntA$ carries its $j$-th channel at $\pslot{\pntA}{j}$.
Every point of a judgment is built from a point variable or a constant.
$\pvivec$ stands for a sequence $\pviA_1 \dots \pviA_m$ of point
variables and $\pntvec$ for a sequence of points, and
$\tysub{\ty}{\pntvec}{\pvivec}$ replaces each $\pviA_i$ by $\pntA_i$.

\begin{figure}[tp]
  \adjustfigure[]
  \rulesection*{Types $\ty$, session types $\sty$, bounds $\ubd, \lbd$ (\rocq{ty})}
  \begin{align*}
    \ty,\tyB &::= \tyUnit \sepr \lPair{\ty}{\tyB} \sepr \lSum{\ty}{\tyB}
      \sepr \ty \lolli{\ubd}{\lbd} \tyB \sepr \sAt{\sty}{\pntA}
      \sepr \tyBox{\pvivec}{\cstr}{\ty} \sepr \tyEx{\lbd}{\pvivec}{\cstr}{\ty}
    \\
    \sty &::= \sSend{\ty}{\sty} \sepr \sRecv{\ty}{\sty} \sepr \sTerm
      \sepr \tyMu\tvar.\sty \sepr \tvar \sepr \dualv{\tvar}
    \\
    \ubd &::= \pBot \sepr \pntA
    & \lbd &::= \pTop \sepr \pBot \sepr \pntA
  \end{align*}
  \rulesection{Session duality $\dual{\sty}$ (\rocq{dual})}
  \begin{align*}
    \dual{\sSend{\ty}{\sty}} &= \sRecv{\ty}{\dual{\sty}}
    &
    \dual{\sRecv{\ty}{\sty}} &= \sSend{\ty}{\dual{\sty}}
    &
    \dual{\sTerm} &= \sTerm
    \\
    \dual{\tvar} &= \dualv{\tvar}
    &
    \dual{\dualv{\tvar}} &= \tvar
    &
    \dual{\tyMu\tvar.\sty} &= \tyMu\tvar.\,\sflipop_{\tvar}(\dual{\sty})
  \end{align*}
  \rulesection{Point substitution $\tysub{\pntA}{\pntvec}{\pvivec}$, $\tysub{\ty}{\pntvec}{\pvivec}$ (\rocq{pt_inst}, \rocq{pinst_ty})}
  \begin{align*}
    \tysub{(\poff{\pviA_i}{k})}{\pntvec}{\pvivec} &= \poff{\pntA_i}{k}
    &
    \tysub{(\poff{\pcst}{k})}{\pntvec}{\pvivec} &= \poff{\pcst}{k}
    &
    \tysub{(\poff{\pslot{\pntB}{j}}{k})}{\pntvec}{\pvivec}
      &= \poff{\pslot{(\tysub{\pntB}{\pntvec}{\pvivec})}{j}}{k}
    \\
    \tysub{(\sAt{\sty}{\pntB})}{\pntvec}{\pvivec}
      &= \sAt{\sty}{(\tysub{\pntB}{\pntvec}{\pvivec})}
    \span\span\span\span
  \end{align*}
  and structurally on the other type formers, mapping over the bounds
  $\ubd$ and $\lbd$ of an arrow and of an existential, and going under
  the binder of a $\Box$ and of an $\exists$ (\rocq{pup}).
  \rulesection{Payloads (\rocq{payload_at})}
  \begin{align*}
    \tysub{\ty}{\pntA}{\pself} &= \tysub{\ty}{\pslot{\pntA}{j}}{\pslot{\pself}{j}}
  \end{align*}
  \rulesectionend
  \caption{Types, session types and duality}
  \Description{The grammars of types, session types and the two kinds of
    bound; the six equations of session duality; and the definition of a
    payload instantiated at a point.}
  \label{fig:types}
\end{figure}

A channel of type $\sAt{\sty}{\pntA}$ runs $\sty$ and acts next at
$\pntA$. The arrow $\ty \lolli{\ubd}{\lbd} \tyB$ carries the latest
point at which a call may block, $\pBot$ for never, and a lower bound on
the points of the channels the closure holds, $\pTop$ for none. Arrows
take $\lbd \in \set{\pTop} \cup \{\pntA\}$ and only existentials take
$\lbd = \pBot$. $\tyBox{\pvivec}{\cstr}{\ty}$ is replicable, holds no
channel, and is usable at every choice of points for $\pvivec$ that
satisfies $\cstr$. $\tyEx{\lbd}{\pvivec}{\cstr}{\ty}$ is $\ty$ at some
such choice.

\begin{definition}[Payload \rocq{payload_at}]
\label{def:payload}
  The type $\ty$ in $\sSend{\ty}{\sty}$ and $\sRecv{\ty}{\sty}$ mentions
  no point other than the markers $\pslot{\pself}{0}, \pslot{\pself}{1},
  \dots$, and a send or receive at $\pntA$ fills $\pslot{\pself}{j}$
  with $\pslot{\pntA}{j}$, written $\tysub{\ty}{\pntA}{\pself}$.
\end{definition}

\begin{remark}
\label{rm:payload}
  A marker of a payload is a self-marker by construction, so no rule
  enforces \cref{def:payload}. Point substitution replaces the
  variables $\pvivec$, and $\pslot{\pself}{j}$ is no variable, so no
  enclosing $\Box$ or $\exists$ touches it. Point substitution also
  stops at a channel type, by the clause
  $\tysub{(\sAt{\sty}{\pntB})}{\pntvec}{\pvivec} =
    \sAt{\sty}{(\tysub{\pntB}{\pntvec}{\pvivec})}$ of \cref{fig:types},
  so it never reaches inside a payload at all (\rocq{pinst_ty}).

  Session substitution does reach inside a payload, because unfolding
  $\tyMu\tvar.\sRecv{(\sAt{\tvar}{\pslot{\pself}{0}})}{\sTerm}$ has to
  rewrite the $\tvar$ under the carrier. It rewrites session variables
  and nothing else, and the point a channel type hangs on stays put
  (\rocq{sinst_ty}, and the same shape in \rocq{sren_ty} and
  \rocq{sflip_ty}), so unfolding leaves every marker alone. The
  property the proofs use is
  $\tysub{(\tysub{\ty}{\pntA}{\pself})}{\pntvec}{\pvivec}
    = \tysub{\ty}{(\tysub{\pntA}{\pntvec}{\pvivec})}{\pself}$
  (\rocq{payload_pinst}), which holds with no side condition because
  $\tysub{\ty}{\pntA}{\pself}$ is a total substitution.
\end{remark}

Sender and receiver name the carried channels alike, so a message
carries no point. $\dualv{\tvar}$ is the dual polarity of $\tvar$, the
dual of whatever $\tvar$ stands for. Source types never contain
$\dualv{\tvar}$: it arises only from duality, and dualizing
$\tyMu\tvar.\sty$ swaps the polarities of its own variable so that
unfolding commutes with duality (\cref{sec:app-duality}).

\paragraph{Contexts and the judgment}
A linear environment $\lenv ::= \tEmpty \sepr \lenv, \vr : \ty$ is
linear except at $\Box$ types. $\lenv[\vr : \ty]$ updates $\lenv$, which
already binds $\vr$, so that $\vr$ has type $\ty$, and it extends
nothing. The facts $\fctx$ assumed by a judgment
are facts $\fact$ over points together with facts $\fLt{\ubd}{\pfv}$ for
frame variables $\pfv$. Nodes are
$\pnode ::= \pntA \sepr \pfv \sepr \pBot$, and a frame $\frm$ is a
finite set of nodes. The judgment
\[
  \fctx; \lenv \typesF{\frm} \tm : \ty
  \qquad(\rocq{styped})
\]
says that $\tm$ uses exactly $\lenv$ and has type $\ty$, and that the
rest of the program holds the nodes of $\frm$ while $\tm$ runs. A frame
variable $\pfv$ stands for whatever the caller of a function holds, and
$\pBot$ for channels at points not known here.

\begin{figure}[tp]
  \adjustfigure[]
  \rulesection*{Aliasing $\fctx \types \fEq{\pntA}{\pntB}$ (\rocq{seqs})}
  \begin{proofrules}
    \infer*[lab=A-Hyp]{
      (\fEq{\pntA}{\pntB}) \in \fctx
    }{
      \fctx \types \fEq{\pntA}{\pntB}
    } \label{rule:A-Hyp}
    \and
    \infer*[lab=A-Refl]{}{
      \fctx \types \fEq{\pntA}{\pntA}
    } \label{rule:A-Refl}
    \and
    \infer*[lab=A-Sym]{
      \fctx \types \fEq{\pntA}{\pntB}
    }{
      \fctx \types \fEq{\pntB}{\pntA}
    } \label{rule:A-Sym}
    \and
    \infer*[lab=A-Trans]{
      \fctx \types \fEq{\pntA}{\pntB} \\
      \fctx \types \fEq{\pntB}{\pntC}
    }{
      \fctx \types \fEq{\pntA}{\pntC}
    } \label{rule:A-Trans}
    \and
    \infer*[lab=A-Shift]{
      \fctx \types \fEq{\pntA}{\pntB}
    }{
      \fctx \types \fEq{\poff{\pntA}{1}}{\poff{\pntB}{1}}
    } \label{rule:A-Shift}
    \and
    \infer*[lab=A-Slot]{
      \fctx \types \fEq{\pntA}{\pntB}
    }{
      \fctx \types \fEq{\pslot{\pntA}{j}}{\pslot{\pntB}{j}}
    } \label{rule:A-Slot}
  \end{proofrules}
  \rulesection{Precedence $\fctx \types \fLt{\pnode}{\pnode'}$ (\rocq{fl})}
  \begin{proofrules}
    \infer*[lab=O-Po]{}{
      \fctx \types \fLt{\pntA}{\poff{\pntA}{1}}
    } \label{rule:O-Po}
    \and
    \infer*[lab=O-Hyp]{
      (\fLt{\pnode}{\pnode'}) \in \fctx
    }{
      \fctx \types \fLt{\pnode}{\pnode'}
    } \label{rule:O-Hyp}
    \and
    \infer*[lab=O-EqL]{
      \fctx \types \fEq{\pntA}{\pntA'} \\
      \fctx \types \fLt{\pntA'}{\pnode}
    }{
      \fctx \types \fLt{\pntA}{\pnode}
    } \label{rule:O-EqL}
    \and
    \infer*[lab=O-EqR]{
      \fctx \types \fLt{\pnode}{\pntB'} \\
      \fctx \types \fEq{\pntB'}{\pntB}
    }{
      \fctx \types \fLt{\pnode}{\pntB}
    } \label{rule:O-EqR}
    \and
    \infer*[lab=O-Trans]{
      \fctx \types \fLt{\pnode}{\pntB} \\
      \fctx \types \fLt{\pntB}{\pnode'}
    }{
      \fctx \types \fLt{\pnode}{\pnode'}
    } \label{rule:O-Trans}
  \end{proofrules}
  \rulesection{Heads, lower bounds, precedence of a frame (\rocq{heads})}
  \begin{align*}
    \heads(\sAt{\sty}{\pntA}) &= \set{\pntA}
    &
    \heads(\ty \lolli{\ubd}{\lbd} \tyB) &= \heads(\tyEx{\lbd}{\pvivec}{\cstr}{\ty}) = \set{\lbd}
    \\
    \heads(\lPair{\ty}{\tyB}) &= \heads(\lSum{\ty}{\tyB}) = \heads(\ty) \cup \heads(\tyB)
    \\
    \heads(\tyUnit) &= \heads(\tyBox{\pvivec}{\cstr}{\ty}) = \emptyset
  \end{align*}
  with $\pTop$ contributing nothing, and
  $\heads(\lenv) = \bigcup \set{\heads(\ty) \mid \vr : \ty \in \lenv}$.
  \begin{align*}
    \fctx \types \fLe{\lbd}{N}
    &\quad\text{iff}\quad
    \lbd = \pTop \text{ and } N = \emptyset,
    \text{ or } \lbd = \pBot,
    \\ & \qquad\qquad
    \text{ or } \lbd = \pntA \text{ and }
      \fctx \types \fEq{\pntA}{\pnode} \text{ or } \fctx \types \fLt{\pntA}{\pnode}
      \text{ for every } \pnode \in N
    \\
    \fctx \types \fLt{\ubd}{\frm}
    &\quad\text{iff}\quad
    \ubd = \pBot,
    \text{ or } \ubd = \pntA \text{ and }
      \fctx \types \fLt{\pntA}{\pnode} \text{ for every } \pnode \in \frm
  \end{align*}
  \rulesectionend
  \caption{The fact theory}
  \Description{Six rules for aliasing, five rules for precedence, the
    heads of a type and of an environment, and the two derived
    judgments relating a bound to a set of nodes.}
  \label{fig:facts}
\end{figure}

\Cref{rule:O-Po} is program order, and it relates two points of one name
only. No rule derives $\fLt{\pnode}{\pBot}$, so a $\pBot$ node in $N$ or
in $\frm$ makes both derived judgments fail for a point: no closure,
message, packed value or blocking call coexists with an unopened
existential at $\pBot$.

\paragraph{Term rules}
\Cref{fig:typeSys,fig:typeSysPoints} give the term rules. $\lenv_1 \cdot \lenv_2$ splits
$\lenv$ and copies variables of $\Box$ type to both sides. The first
subterm of every composite runs while the rest holds its channels, so it
is typed under $\frm$ grown by the heads of that rest.
\ref{rule:E-Abs} types the body under the frame $\set{\pfv}$ and assumes
$\fLt{\ubd}{\pfv}$, so every call in the body precedes $\ubd$, and
\ref{rule:E-App} checks that $\ubd$ precedes what the caller holds. The
two checks meet at $\ubd$.
\ref{rule:E-Exp} is the one recursive rule and the only rule that binds
$\forall$: $\forall$ appears only directly under $\Box$, because a value
that holds channels has fixed points.

\begin{figure}[tp]
  \adjustfigure[]
  \rulesection*{Typing rules for terms $\fctx ; \lenv \typesF{\frm} \tm : \ty$ (\rocq{styped})}
  \begin{proofrules}
    \infer*[lab=E-Unit]{}{
      \fctx ; \tEmpty \typesF{\frm} \tUnit : \tyUnit
    } \label{rule:E-Unit}
    \and
    \infer*[lab=E-Const]{
      \cst : \ty \text{ an instance of a signature under } \fctx
    }{
      \fctx ; \tEmpty \typesF{\frm} \cst : \ty
    } \label{rule:E-Const}
    \and
    \infer*[lab=E-Var]{}{
      \fctx ; \vr : \ty \typesF{\frm} \vr : \ty
    } \label{rule:E-Var}
    \and
    \infer*[lab=E-Abs]{
      \pfv \text{ fresh} \\
      \fctx \types \fLe{\lbd}{\heads(\lenv)} \\
      \fctx, \fLt{\ubd}{\pfv} ; \lenv , \vr : \ty \typesF{\set{\pfv}} \tm : \tyB
    }{
      \fctx ; \lenv \typesF{\frm} \tAbs{\vr}{\tm} : \ty \lolli{\ubd}{\lbd} \tyB
    } \label{rule:E-Abs}
    \and
    \infer*[lab=E-App]{
      \fctx \types \fLt{\ubd}{\frm} \\
      \fctx ; \lenv_1 \typesF{\frm \cup \heads(\lenv_2)} \tm_1 : \ty \lolli{\ubd}{\lbd} \tyB \\
      \fctx ; \lenv_2 \typesF{\frm \cup \set{\lbd}} \tm_2 : \ty
    }{
      \fctx ; \lenv_1 \cdot \lenv_2 \typesF{\frm} \tApp{\tm_1}{\tm_2} : \tyB
    } \label{rule:E-App}
    \and
    \infer*[lab=E-Pair]{
      \fctx ; \lenv_1 \typesF{\frm \cup \heads(\lenv_2)} \tm_1 : \ty \\
      \fctx ; \lenv_2 \typesF{\frm \cup \heads(\ty)} \tm_2 : \tyB
    }{
      \fctx ; \lenv_1 \cdot \lenv_2 \typesF{\frm} \tPair{\tm_1}{\tm_2} : \lPair{\ty}{\tyB}
    } \label{rule:E-Pair}
    \and
    \infer*[lab=E-LetPair]{
      \fctx ; \lenv_1 \typesF{\frm \cup \heads(\lenv_2)} \tm_1 : \lPair{\tyB}{\tyC} \\
      \fctx ; \lenv_2 , \vr : \tyB , \vrB : \tyC \typesF{\frm} \tm_2 : \ty
    }{
      \fctx ; \lenv_1 \cdot \lenv_2 \typesF{\frm} \tLet[\vr][\vrB]{\tm_1}{\tm_2} : \ty
    } \label{rule:E-LetPair}
    \and
    \infer*[lab=E-Let*]{
      \fctx ; \lenv_1 \typesF{\frm \cup \heads(\lenv_2)} \tm_1 : \tyUnit \\
      \fctx ; \lenv_2 \typesF{\frm} \tm_2 : \ty
    }{
      \fctx ; \lenv_1 \cdot \lenv_2 \typesF{\frm} \tLet*{\tm_1}{\tm_2} : \ty
    } \label{rule:E-Let*}
    \and
    \infer*[lab=E-Inj]{
      \fctx ; \lenv \typesF{\frm} \tm : \ty_{\asTy{i}}
    }{
      \fctx ; \lenv \typesF{\frm} \tInj{\asTm{i}}{\tm} : \lSum{\ty_1}{\ty_2}
    } \label{rule:E-Inj}
    \and
    \infer*[lab=E-Case]{
      \fctx ; \lenv_1 \typesF{\frm \cup \heads(\lenv_2)} \tm : \lSum{\ty_1}{\ty_2} \\
      \fctx ; \lenv_2 , \ivx_1 : \ty_1 \typesF{\frm} \tm_1 : \tyB \\
      \fctx ; \lenv_2 , \ivx_2 : \ty_2 \typesF{\frm} \tm_2 : \tyB
    }{
      \fctx ; \lenv_1 \cdot \lenv_2 \typesF{\frm} \tCase{\tm}{\ivx_1}{\tm_1}{\ivx_2}{\tm_2} : \tyB
    } \label{rule:E-Case}
  \end{proofrules}
  \rulesectionend
  \caption{Typing rules for terms: variables, functions and data}
  \Description{Ten term typing rules, for the unit, constants,
    variables, abstraction, application, pairs, their elimination,
    sequencing, injections and case.}
  \label{fig:typeSys}
\end{figure}

\begin{figure}[tp]
  \adjustfigure[]
  \rulesection*{Structural rules and rules for points, $\fctx ; \lenv \typesF{\frm} \tm : \ty$ (\rocq{styped})}
  \begin{proofrules}
    \infer*[lab=E-Weak]{
      \fctx ; \lenv \typesF{\frm} \tm : \ty \\
      \lenv' \text{ only } \Box \text{ types}
    }{
      \fctx ; \lenv , \lenv' \typesF{\frm} \tm : \ty
    } \label{rule:E-Weak}
    \and
    \infer*[lab=E-Contract]{
      \fctx ; \lenv , \ivx : \tBox{\ty} , \ivy : \tBox{\ty} \typesF{\frm} \tm : \tyB
    }{
      \fctx ; \lenv , \ivx : \tBox{\ty} \typesF{\frm} \tSub{\tm}{\ivx}{\ivy} : \tyB
    } \label{rule:E-Contract}
    \and
    \infer*[lab=E-Conv]{
      \fctx ; \lenv \typesF{\frm} \tm : \ty \\
      \ty \approx \ty' \text{ up to } \fctx \types {\equiv}
    }{
      \fctx ; \lenv \typesF{\frm} \tm : \ty'
    } \label{rule:E-Conv}
    \and
    \infer*[lab=E-Exp]{
      \lenv \text{ only } \Box \text{ types} \\
      \pvivec, \pfv \text{ fresh} \\\\
      \fctx, \cstr, \fLt{\ubd}{\pfv} ;
        \lenv , \vr : \tyBox{\pvivec}{\cstr}{\tyUnit \lolli{\ubd}{\pTop} \ty}
        \typesF{\set{\pfv}} \tm : \ty
    }{
      \fctx ; \lenv \typesF{\frm} \tBoxP{\vr}{\tm}
        : \tyBox{\pvivec}{\cstr}{\tyUnit \lolli{\ubd}{\pTop} \ty}
    } \label{rule:E-Exp}
    \and
    \infer*[lab=E-Pack]{
      \fctx \types \tysub{\cstr}{\pntvec}{\pvivec} \\
      \fctx \types \fLe{\lbd}{\heads(\tysub{\ty}{\pntvec}{\pvivec})} \\
      \fctx ; \lenv \typesF{\frm} \tm : \tysub{\ty}{\pntvec}{\pvivec}
    }{
      \fctx ; \lenv \typesF{\frm} \tm : \tyEx{\lbd}{\pvivec}{\cstr}{\ty}
    } \label{rule:E-Pack}
    \and
    \infer*[lab=E-Open]{
      \vr : \tyEx{\lbd}{\pvivec}{\cstr}{\ty} \in \lenv \\
      \pvivec \text{ fresh, not in } \tyB \\
      \fctx, \cstr ; \lenv[\vr : \ty] \typesF{\frm} \tm : \tyB
    }{
      \fctx ; \lenv \typesF{\frm} \tm : \tyB
    } \label{rule:E-Open}
  \end{proofrules}
  \rulesectionend
  \caption{Typing rules for terms: structure and points}
  \Description{Six term typing rules, for weakening, contraction,
    conversion, a box, packing and opening an existential.}
  \label{fig:typeSysPoints}
\end{figure}

\paragraph{Signatures}
\Cref{fig:signatures} gives the type-indexed family of signatures the
constants are instantiated from. An action that may block happens at a
point $\pntA$, so $\tSend$, $\tRecv$ and $\tClos$ have upper bound
$\pntA$, and \ref{rule:E-App} then requires
$\fctx \types \fLt{\pntA}{\frm}$. $\tLink$ and $\tSpawn$ never block, so
their upper bound is $\pBot$. $\tNew$ puts both ends of a channel at one
point $\pviA$ and returns an existential, so a channel starts at a point
placed by its creation constraint and not by a global number.
$\tLink$ asks $\fEq{\pntA}{\pntB}$ because forwarding merges two
sessions into one conduit, and the actions of the merged channel need a
single position. Under $\fLt{\pntA}{\pntB}$ they would have two. The
condition is not an artefact of the static rule, \ref{rule:O-Link} of
the logical relation requires the world to alias the two points as well.
Whether a weaker condition that coerces one position into the other
would work is open.

\begin{figure}[tp]
  \adjustfigure[]
  \rulesection*{Type-indexed families of type signatures for constants $\cst : \ty$ (\rocq{konst_ty})}
  \begin{align*}
    & & & \text{side conditions the instance discharges}
    \\
    \tNew &{}: \tyUnit \lolli{\pBot}{\pTop}
      \tyEx{\pBot}{\pviA}{\cstr}{\lPair{\sAt{\sty}{\pviA}}{\sAt{\dual{\sty}}{\pviA}}}
    & & \cstr \text{ a creation constraint}
    \\
    \tSpawn &{}: (\tyUnit \lolli{\ubd}{\lbd} \tyUnit) \lolli{\pBot}{\pTop} \tyUnit
    \\
    \tSend &{}: \lPair{\tysub{\ty}{\pntA}{\pself}}{\sAt{(\sSend{\ty}{\sty})}{\pntA}}
      \lolli{\pntA}{\pTop} \sAt{\sty}{(\poff{\pntA}{1})}
    & & \fctx \types \fLt{\pntA}{\pnode} \text{ for every }
        \pnode \in \heads(\tysub{\ty}{\pntA}{\pself})
    \\
    \tRecv &{}: \sAt{(\sRecv{\ty}{\sty})}{\pntA}
      \lolli{\pntA}{\pTop} \lPair{\tysub{\ty}{\pntA}{\pself}}{\sAt{\sty}{(\poff{\pntA}{1})}}
    \\
    \tClos &{}: \sAt{\sTerm}{\pntA} \lolli{\pntA}{\pTop} \tyUnit
    \\
    \tLink &{}: \lPair{\sAt{\sty}{\pntA}}{\sAt{\dual{\sty}}{\pntB}}
      \lolli{\pBot}{\pTop} \tyUnit
    & & \fctx \types \fEq{\pntA}{\pntB}
    \\
    \tForce &{}: \tyBox{\pvivec}{\cstr}{\tyUnit \lolli{\ubd}{\pTop} \ty}
      \lolli{\tysub{\ubd}{\pntvec}{\pvivec}}{\pTop} \tysub{\ty}{\pntvec}{\pvivec}
    & & \fctx \types \tysub{\cstr}{\pntvec}{\pvivec}
  \end{align*}
  \rulesectionend
  \caption{Signatures of the constants}
  \Description{The seven type signatures, with the side conditions on
    points that each one reads from the fact theory.}
  \label{fig:signatures}
\end{figure}

\begin{definition}[Creation constraint \rocq{cc_ok}]
\label{def:creation}
  A point of $\pviA$ is a point with root $\pviA$. A constraint $\cstr$
  over one bound variable $\pviA$ is a \emph{creation constraint} under
  $\fctx$ when it is
  \begin{itemize}
    \item \emph{acyclic}:
      $\fctx, \cstr \nvdash \fLt{\pntA}{\pntA}$ for every point $\pntA$;
      and
    \item \emph{conservative}:
      $\fctx, \cstr \types \fLt{\pntA}{\pntB}$ implies
      $\fctx \types \fLt{\pntA}{\pntB}$, and
      $\fctx, \cstr \types \fEq{\pntA}{\pntB}$ implies
      $\fctx \types \fEq{\pntA}{\pntB}$, for all points
      $\pntA, \pntB$ without root $\pviA$.
  \end{itemize}
\end{definition}

A new channel may sit anywhere the context leaves room, and it may not
order or identify two points the context already has.
$\cstr = \set{\fEq{\pviA}{\pntA}}$ places the channel at $\pntA$, and
$\cstr = \set{\fLt{\ubd}{\pviA}, \fLt{\pviA}{\asTy{v}}}$ places it
between $\ubd$ and $\asTy{v}$, which conservativity allows only when
$\fctx \types \fLt{\ubd}{\asTy{v}}$. We write
$\tNewC{\pviA}{\cstr}$ for $\tNew$ at a fresh point $\pviA$ under
$\cstr$, $\tBoxA{\vr}{\pvivec}{\cstr}{\tm}$ for the box of
\ref{rule:E-Exp}, $\tAbsA{\ubd}{\lbd}{\vr}{\tm}$ for the abstraction of
\ref{rule:E-Abs}, and $\tForceA{\vr}{\pntvec}$ for the instantiation of
\ref{rule:E-Const} at $\tForce$. The derived form
$\tLetS{\tPair{\ch}{\chB}}{\tApp{\tNew}{\tUnit}}{\tm}$ unfolds to
\ref{rule:E-LetPair}, \ref{rule:E-Pair}, \ref{rule:E-App},
\ref{rule:E-Let*}, \ref{rule:E-Open} and \ref{rule:E-LetPair}
(\rocq{t_new_open}).

\paragraph{Processes}
\Cref{fig:procTypeSys} types processes. All return types are
\emph{primitive} (\rocq{prim}), so a thread cannot avoid closing its
channels by returning them. \ref{rule:P-Connect} composes two processes
that interact in a dual way on the channels of $\senv_{\asTy{C}}$ and
subsumes the mix rule at $\senv_{\asTy{C}} = \tEmpty$.

The process judgment $\senv \typesW{\wrld} \proc : \ty$ carries a world
$\wrld$, and so does the term judgment
$\tenv ; \lenv \typesFW{\frm}{\wrld} \tm : \ty$ that
\ref{rule:P-Term} reads (\rocq{typed}). Neither is the judgment
$\fctx ; \lenv \typesF{\frm} \tm : \ty$ of
\cref{fig:typeSys,fig:typeSysPoints}, which carries facts and no world.
\Cref{sec:worlds} defines worlds, \cref{sec:logicalRelation:relation}
defines the term environment $\tenv$, and \rocq{styped_typed} embeds
every source derivation into the world-carrying judgment at the empty
world. The restricted channels of \ref{rule:P-Term} are the domain of
$\tenv_{\asTy{B}}$, which holds both ends of each of them, so a thread
that owns a whole channel is typed without the channel appearing in its
interface.

\begin{figure}[tp]
  \adjustfigure[]
  \rulesection*{Primitive types $\primitive \ty$ (\rocq{prim})}
  \begin{proofrules}
    \infer*{}{
      \primitive \tyUnit
    }
    \and
    \infer*{
      \primitive \tyB \\ \primitive \tyC
    }{
      \primitive \lPair{\tyB}{\tyC}
    }
    \and
    \infer*{
      \primitive \tyB \\ \primitive \tyC
    }{
      \primitive \lSum{\tyB}{\tyC}
    }
    \and
    \infer*{
      \primitive \tyB
    }{
      \primitive \tyBox{\pvivec}{\cstr}{\tyUnit \lolli{\ubd}{\pTop} \tyB}
    }
  \end{proofrules}
  \rulesection{Typing rules for processes $\senv \typesW{\wrld} \proc : \ty$ (\rocq{ptyped})}
  \begin{proofrules}
    \infer*[lab=P-Term]{
      \primitive \ty \\
      \tenv_{\asTy{B}} \text{ holds both ends of each of its channels,
        dual at aliased points} \\
      \senv , \tenv_{\asTy{B}} ; \tEmpty \typesFW{\emptyset}{\wrld} \tm : \ty
    }{
      \senv \typesW{\wrld} (\dom(\tenv_{\asTy{B}}) , \mset{\tm}) : \ty
    } \label{rule:P-Term}
    \and
    \infer*[lab=P-Connect]{
      \senv_1 , \senv_{\asTy{C}} \typesW{\wrld} (\chset_1 , \pool_1) : \ty \\
      \senv_2 , \dual{\senv_{\asTy{C}}} \typesW{\wrld} (\chset_2 , \pool_2) : \tyUnit \\
      \dom(\senv_{\asTy{C}}) = \chset_0
    }{
      \senv_1 , \senv_2 \typesW{\wrld} (\chset_1 \cupdot \chset_2 \cupdot \chset_0 , \pool_1 \uplus \pool_2) : \ty
    } \label{rule:P-Connect}
  \end{proofrules}
  \rulesectionend
  \caption{Typing for processes}
  \Description{The four rules defining primitive types and the two
    process typing rules.}
  \label{fig:procTypeSys}
\end{figure}

\paragraph{Soundness}
A program starts under one constraint $\cstr_0$ over the top-level point
variables it begins with. A program that creates all its channels with
$\tNew$ takes $\cstr_0 = \emptyset$. A program that wants a global
order, such as the translation of a priority system, states it in
$\cstr_0$ and aliases its channels to the variables $\cstr_0$ speaks of.
The source programs of this paper name no point constant. Their points
are the variables that $\cstr_0$, a $\Box$ and an $\exists$ bind, and
\cref{thm:soundness-source} supplies the constants itself
(\rocq{c0_consts}). For such a program $\tysub{\ty}{\pntA}{\pself}$ is
the only way a payload acquires a point. A translated program may name
constants, and the image of PGV does (\rocq{proot}).

\begin{theorem}[Soundness for source programs \rocq{soundness_source}]
\label{thm:soundness-source}
  If $\cstr_0$ is acyclic, $\primitive \ty$ and
  $\cstr_0 ; \tEmpty \typesF{\emptyset} \tm : \ty$, then the thread
  $(\emptyset, \mset{\tm})$ is deadlock free.
\end{theorem}

\begin{proof}
  The proof builds the world the run needs out of $\cstr_0$. Embed the
  source derivation into the world-carrying judgment at the empty world,
  which supports no fact beyond equality and program order
  (\rocq{wnil}, \rocq{wnil_strict}, \rocq{styped_typed}). Let $n$ be the
  number of top-level point variables and send each of them to its own
  point constant. Acyclicity of $\cstr_0$ makes the world generated by
  the image of $\cstr_0$ under that map strict (\rocq{wc0},
  \rocq{wc0_strict}), it extends the empty world (\rocq{wc0_ext}), and
  the hypothesis $(\cstr_0, \pvivec)$ holds in it at those constants
  (\rocq{wc0_holds}). Instantiating the derivation along the map
  therefore yields the infinitary judgment in that world
  (\rocq{typed_sem_rules}, \rocq{c0_consts}), hence a semantic judgment
  (\rocq{sem_rules_fundamental}). One thread with nothing restricted is
  a process of the relation (\rocq{Prel_term}), and
  \cref{thm:adequacy} concludes. Point instantiation preserves
  $\primitive$ (\rocq{prim_pinst}), which \rocq{Prel_term} needs.
\end{proof}

Acyclicity of $\cstr_0$ is the acyclicity half of
\cref{def:creation}, checked against the empty context. The judgment of
\cref{thm:soundness-source} carries no world and no channel environment.
The world-carrying judgment $\senv \typesW{\wrld} \proc : \ty$ of
\cref{fig:procTypeSys} and the logical relation of
\cref{sec:logicalRelation} sit behind it. \rocq{styped_typed} bridges
them, taking a source derivation to the world-carrying term judgment at
any world and at the empty term environment, and \rocq{typed_styped}
comes back from a world with empty support.

\begin{example}[Delegation \rocq{deleg_main}]
\label{ex:delegation}
  One channel carries another, and the delegated session runs strictly
  inside the gap between the carrier's two actions. The two protocols
  are
  \begin{align*}
    \sty^0 &= \sSend{\tyUnit}{\sTerm}
    & \sty &= \sSend{(\sAt{\sty^0}{\pslot{\pself}{0}})}{\sTerm}
  \end{align*}
  and the program, with every annotation inlined, is
  \begin{align*}
    \asTmId{deleg} := {}
      & \tLetS{\tPair{\ch : \sAt{\sty}{\pviA}}{\ch' : \sAt{\dual{\sty}}{\pviA}}}{%
          \tApp{\tNewC{\pviA}{\fLt{\pviA}{\pslot{\pviA}{0}},\ %
            \fLt{\poff{\pslot{\pviA}{0}}{1}}{\poff{\pviA}{1}}}}{\tUnit}}{}
      \\
      & \tLetS{\tPair{\chB : \sAt{\sty^0}{\pslot{\pviA}{0}}}{%
          \chB' : \sAt{(\sRecv{\tyUnit}{\sTerm})}{\pslot{\pviA}{0}}}}{%
          \tApp{\tNewC{\pviB}{\fEq{\pviB}{\pslot{\pviA}{0}}}}{\tUnit}}{}
      \\
      & \tSeq{\tApp{\tSpawn}{\asTmId{K}}}{}
        \tLetS{\ch_1 : \sAt{\sTerm}{(\poff{\pviA}{1})}}{\tApp{\tSend}{\tPair{\chB}{\ch}}}{}
      \\
      & \tLet[\ivx][\chB_1]{\tApp{\tRecv}{\chB'}}{}
        \tSeq{\ivx}{\tSeq{\tApp{\tClos}{\chB_1}}{\tApp{\tClos}{\ch_1}}}
      \\
    \asTmId{K} := {}
      & \tAbsA{\poff{\pviA}{1}}{\pviA}{\tUnit}{}
        \tLet[\asTm{e}][\ch'_1]{\tApp{\tRecv}{\ch'}}{}
        \tSeq{\tApp{\tClos}{(\tApp{\tSend}{\tPair{\tUnit}{\asTm{e}}})}}{\tApp{\tClos}{\ch'_1}}
  \end{align*}
  The creation constraint of $\ch$ puts the delegated session strictly
  inside the gap between the carrier's two actions, and the second
  $\tNew$ creates $\chB$ at the carried root of a message the carrier
  has not sent yet. The order is one chain,
  \[
    \fLt{\pviA}{{}} \pslot{\pviA}{0}
    \prec \poff{\pslot{\pviA}{0}}{1}
    \prec \poff{\pviA}{1}
    \prec \pfv ,
  \]
  built from two facts of the first $\tNew$, program order, and the
  frame hypothesis \ref{rule:E-Abs} adds for the spawned thread. The
  rank $\poff{\pviA}{k} \mapsto 10k$ and
  $\poff{\pslot{\pviA}{0}}{o} \mapsto 5+o$ witnesses acyclicity.
  \Cref{sec:app-delegation} lists every blocking action with the premise
  it owes. The judgment is
  $\emptyset ; \tEmpty \typesF{\emptyset} \asTmId{deleg} : \tyUnit$
  (\rocq{deleg_typed}), so $\asTmId{deleg}$ is deadlock free
  (\rocq{deleg_deadlock_free}).

  A priority system cannot type this program. The carrier typed
  $\sSend{}{}^0.\sRecv{}{}^1$ leaves no integer between its two actions,
  and the number of actions the delegated session performs is not known
  where the carrier is typed.
\end{example}

%%% Local Variables:
%%% mode: latex
%%% TeX-master: "main"
%%% End:
